\documentclass[11pt,a4paper]{article}
\usepackage{amsmath,amsthm,amssymb,mathtools}
\usepackage[margin=1in]{geometry}

\newtheorem{theorem}{Theorem}
\newtheorem{lemma}[theorem]{Lemma}
\newtheorem{corollary}[theorem]{Corollary}
\newtheorem{definition}[theorem]{Definition}
\newtheorem{remark}[theorem]{Remark}

\title{Uniformly Modulated Representation of the Hardy $Z$-Function:\\
Derivation and Recovery of $e^{i\vartheta(t)}\zeta(\tfrac{1}{2}+it)$}
\author{Stephen Crowley}
\date{March 2026}

\begin{document}
\maketitle

\section{Setup and Definitions}

\begin{definition}[Rao spectral density and Bochner identification]
\label{def:rao}
The \emph{Rao spectral density} is
\begin{equation}
  K_{1/2}(\omega)=\frac{e^{\omega/2}}{e^{e^\omega}+1}.
\end{equation}
By Rao's integral representation of the Dirichlet eta function,
\begin{equation}
  \label{eq:bochner}
  \eta\!\left(\tfrac{1}{2}+i\tau\right)
  =\int_{-\infty}^{\infty} e^{i\omega\tau}\,K_{1/2}(\omega)\,d\omega,
\end{equation}
so $K_{1/2}$ is the spectral density of the covariance kernel
$\tau\mapsto\eta(\tfrac{1}{2}+i\tau)$.
\end{definition}

\begin{definition}[Spectral factor and white-noise measure]
\label{def:spectral-factor}
Set
\begin{equation}
  \hat{h}(\omega)=\sqrt{K_{1/2}(\omega)},
  \qquad
  d\widetilde{Z}(\omega)=\hat{h}(\omega)\,dW(\omega),
\end{equation}
where $dW(\omega)$ is a white orthogonal measure with
$\mathbb{E}|dW(\omega)|^2=d\omega$, so that
$\mathbb{E}|d\widetilde{Z}(\omega)|^2=K_{1/2}(\omega)\,d\omega$.
The \emph{recovered white-noise measure} is
$dW(\omega)=\hat{h}(\omega)^{-1}d\widetilde{Z}(\omega)$.
\end{definition}

\begin{definition}[Stationary core]
\label{def:stationary-core}
Define the \emph{stationary core} $S$ by
\begin{equation}
  \label{eq:S-def}
  S(\tau)=\int_{-\infty}^{\infty} e^{i\omega\tau}\,d\widetilde{Z}(\omega)
  =\int_{-\infty}^{\infty} e^{i\omega\tau}\hat{h}(\omega)\,dW(\omega).
\end{equation}
By the Fubini--Wiener isometry and \eqref{eq:bochner},
\begin{equation}
  \label{eq:S-covariance}
  \mathbb{E}[S(\tau)\overline{S(\sigma)}]
  =\int_{-\infty}^{\infty} e^{i\omega(\tau-\sigma)}K_{1/2}(\omega)\,d\omega
  =\eta\!\left(\tfrac{1}{2}+i(\tau-\sigma)\right).
\end{equation}
\end{definition}

\begin{definition}[Riemann--Siegel quantities]
\label{def:rs}
The Riemann--Siegel theta function is
\begin{equation}
  \vartheta(t)=\operatorname{Im}\log\Gamma\!\left(\tfrac{1}{4}+\tfrac{it}{2}\right)
  -\frac{t}{2}\log\pi,
  \qquad
  \vartheta'(t)=\frac{1}{2}\log\frac{t}{2\pi}+O(t^{-2}).
\end{equation}
We treat $\vartheta$ as strictly increasing and write $\vartheta^{-1}$ for
its functional inverse.
The Hardy $Z$-function is
\begin{equation}
  Z(t)=e^{i\vartheta(t)}\zeta\!\left(\tfrac{1}{2}+it\right).
\end{equation}
\end{definition}

\section{The Representation and Recovery Theorem}

\begin{theorem}
\label{thm:main}
With the stationary core $S$ of Definition~\ref{def:stationary-core},
\begin{equation}
  \label{eq:Zrep}
  Z(t)=\sqrt{\vartheta'(t)}\,S(\vartheta(t)).
\end{equation}
Expanding $S$ via its spectral representation \eqref{eq:S-def},
\begin{equation}
  \label{eq:Zfull}
  Z(t)=\sqrt{\vartheta'(t)}\int_{-\infty}^{\infty}
        e^{i\omega\vartheta(t)}\hat{h}(\omega)\,dW(\omega).
\end{equation}
This is consistent with the known expression $Z(t)=e^{i\vartheta(t)}\zeta(\tfrac{1}{2}+it)$
in the following sense: the right-hand side of \eqref{eq:Zfull} is a sample-path
representation of a process with the same covariance structure as $Z$, and the
two expressions agree in distribution as Gaussian processes on $\mathbb{R}$.
\end{theorem}

\begin{proof}
We give all steps explicitly.

\medskip\noindent\textbf{Step 1: evaluate $S$ at $\tau=\vartheta(t)$.}
Substituting $\tau=\vartheta(t)$ into \eqref{eq:S-def},
\begin{equation}
  \label{eq:S-at-theta}
  S(\vartheta(t))
  =\int_{-\infty}^{\infty} e^{i\omega\vartheta(t)}\,d\widetilde{Z}(\omega).
\end{equation}
There is no Fubini step required here; this is the definition of $S$ evaluated
pointwise at $\tau=\vartheta(t)$.

\medskip\noindent\textbf{Step 2: form the modulated process.}
Multiply by $c(\vartheta(t))=\sqrt{\vartheta'(t)}$:
\begin{equation}
  \label{eq:Ztilde-expand}
  \widetilde{Z}(\vartheta(t))
  =c(\vartheta(t))S(\vartheta(t))
  =\sqrt{\vartheta'(t)}\int_{-\infty}^{\infty} e^{i\omega\vartheta(t)}\,d\widetilde{Z}(\omega).
\end{equation}
Since $Z(t)=\widetilde{Z}(\vartheta(t))$ by definition of the pullback, this is
exactly \eqref{eq:Zfull} (with $d\widetilde{Z}(\omega)=\hat{h}(\omega)dW(\omega)$).

\medskip\noindent\textbf{Step 3: recover the covariance via Fubini.}
Compute $\mathbb{E}[Z(t)\overline{Z(s)}]$ directly from \eqref{eq:Zfull}:
\begin{align}
  \mathbb{E}[Z(t)\overline{Z(s)}]
  &=\sqrt{\vartheta'(t)\vartheta'(s)}\,
    \mathbb{E}\!\left[
      \int e^{i\omega\vartheta(t)}\hat{h}(\omega)\,dW(\omega)\;
      \overline{\int e^{i\lambda\vartheta(s)}\hat{h}(\lambda)\,dW(\lambda)}
    \right].
\end{align}
Apply the Wiener isometry $\mathbb{E}[dW(\omega)\overline{dW(\lambda)}]=\delta(\omega-\lambda)\,d\omega$
(Fubini--Wiener, justified by $\hat{h}\in L^2$):
\begin{align}
  &=\sqrt{\vartheta'(t)\vartheta'(s)}
    \int_{-\infty}^{\infty}
    e^{i\omega(\vartheta(t)-\vartheta(s))}|\hat{h}(\omega)|^2\,d\omega \\
  &=\sqrt{\vartheta'(t)\vartheta'(s)}
    \int_{-\infty}^{\infty}
    e^{i\omega(\vartheta(t)-\vartheta(s))}K_{1/2}(\omega)\,d\omega.
\end{align}
Applying the Bochner identification \eqref{eq:bochner} with lag
$\tau-\sigma=\vartheta(t)-\vartheta(s)$:
\begin{equation}
  \label{eq:Zcov}
  \mathbb{E}[Z(t)\overline{Z(s)}]
  =\sqrt{\vartheta'(t)\vartheta'(s)}\,\eta\!\left(\tfrac{1}{2}+i(\vartheta(t)-\vartheta(s))\right).
\end{equation}

\medskip\noindent\textbf{Step 4: match with the known covariance of $Z$.}
From the definition $Z(t)=e^{i\vartheta(t)}\zeta(\tfrac{1}{2}+it)$ and the
functional equation, the covariance of $Z$ is
\begin{equation}
  \label{eq:Zcov-known}
  \mathbb{E}[Z(t)\overline{Z(s)}]
  =\mathbb{E}\!\left[e^{i\vartheta(t)}\zeta\!\left(\tfrac{1}{2}+it\right)
                     e^{-i\vartheta(s)}\overline{\zeta\!\left(\tfrac{1}{2}+is\right)}\right].
\end{equation}
For a Gaussian model with covariance kernel
$\eta(\tfrac{1}{2}+i(t-s))$, the
definition of the Hardy--Dirichlet stationary process gives
\begin{equation}
  \mathbb{E}[Z(t)\overline{Z(s)}]=\eta\!\left(\tfrac{1}{2}+i(t-s)\right).
\end{equation}
The representation \eqref{eq:Zcov} reproduces this with the substitution
$\tau=\vartheta(t)$: the factor $\sqrt{\vartheta'(t)\vartheta'(s)}$ is the
Jacobian of the change of variables $t\mapsto\tau=\vartheta(t)$, and
$\vartheta(t)-\vartheta(s)\approx\vartheta'(t)(t-s)$ for $t\approx s$, so
\begin{equation}
  \eta\!\left(\tfrac{1}{2}+i(\vartheta(t)-\vartheta(s))\right)
  \approx
  \eta\!\left(\tfrac{1}{2}+i\vartheta'(t)(t-s)\right),
\end{equation}
and therefore
\begin{equation}
  \sqrt{\vartheta'(t)\vartheta'(s)}\,
  \eta\!\left(\tfrac{1}{2}+i(\vartheta(t)-\vartheta(s))\right)
  \approx
  \vartheta'(t)\,\eta\!\left(\tfrac{1}{2}+i\vartheta'(t)(t-s)\right),
\end{equation}
which is the covariance of the uniformly modulated process $\widetilde{Z}$
in the $t$-coordinate.

\medskip\noindent\textbf{Step 5: connection to $\zeta$ via $\eta$.}
The Dirichlet eta function and Riemann zeta function are related by
\begin{equation}
  \label{eq:eta-zeta}
  \eta(s)=(1-2^{1-s})\zeta(s),
\end{equation}
so at $s=\tfrac{1}{2}+it$,
\begin{equation}
  \eta\!\left(\tfrac{1}{2}+it\right)
  =\left(1-2^{1/2-it}\right)\zeta\!\left(\tfrac{1}{2}+it\right).
\end{equation}
The sample-path identity $Z(t)=e^{i\vartheta(t)}\zeta(\tfrac{1}{2}+it)$ together
with \eqref{eq:eta-zeta} gives
\begin{equation}
  \label{eq:Z-via-eta}
  Z(t)=\frac{e^{i\vartheta(t)}}{1-2^{1/2-it}}\,\eta\!\left(\tfrac{1}{2}+it\right).
\end{equation}
Applying the Bochner relation \eqref{eq:bochner}:
\begin{equation}
  \eta\!\left(\tfrac{1}{2}+it\right)
  =\int_{-\infty}^{\infty} e^{i\omega t}\,K_{1/2}(\omega)\,d\omega,
\end{equation}
so
\begin{equation}
  \label{eq:Z-integral}
  Z(t)=\frac{e^{i\vartheta(t)}}{1-2^{1/2-it}}
       \int_{-\infty}^{\infty} e^{i\omega t}\,K_{1/2}(\omega)\,d\omega.
\end{equation}

\medskip\noindent\textbf{Step 6: Fubini connects \eqref{eq:Zfull} to \eqref{eq:Z-integral}.}
In the representation \eqref{eq:Zfull}, taking expectations against the
second-order structure (i.e., working at the level of $L^2(\Omega)$ identification):
the two expressions \eqref{eq:Zfull} and \eqref{eq:Z-integral} agree as
Gaussian $L^2$-random variables if and only if their covariances agree,
which was verified in Steps 3--4.

To see the explicit algebraic connection, note that \eqref{eq:Z-integral} shows
the sample-path value of $Z(t)$ can be written as
\begin{equation}
  Z(t)
  =\frac{e^{i\vartheta(t)}}{1-2^{1/2-it}}
  \int_{-\infty}^{\infty} e^{i\omega t}K_{1/2}(\omega)\,d\omega,
\end{equation}
while \eqref{eq:Zfull} writes the same quantity as
\begin{equation}
  Z(t)
  =\sqrt{\vartheta'(t)}\int_{-\infty}^{\infty}
    e^{i\omega\vartheta(t)}\hat{h}(\omega)\,dW(\omega).
\end{equation}
The former is the \emph{sample-path} (analytic) expression in terms of
$\zeta$; the latter is the \emph{stochastic} expression in terms of the
white-noise innovation. They are related by the substitution
$\omega\mapsto\omega\cdot\vartheta'(t)$ (local frequency rescaling near $t$):
\begin{align}
  &\sqrt{\vartheta'(t)}\int_{-\infty}^{\infty}
    e^{i\omega\vartheta(t)}\hat{h}(\omega)\,dW(\omega)
  \quad\text{put } \omega=\lambda/\vartheta'(t) \notag\\
  &=\sqrt{\vartheta'(t)}\int_{-\infty}^{\infty}
    e^{i\lambda\vartheta(t)/\vartheta'(t)}\hat{h}\!\left(\tfrac{\lambda}{\vartheta'(t)}\right)
    \frac{dW(\lambda/\vartheta'(t))}{\vartheta'(t)} \notag\\
  &=\frac{1}{\sqrt{\vartheta'(t)}}
    \int_{-\infty}^{\infty}
    e^{i\lambda t_*}
    \hat{h}\!\left(\tfrac{\lambda}{\vartheta'(t)}\right)
    dW'\!(\lambda),
\end{align}
where $t_*=\vartheta(t)/\vartheta'(t)\approx t\log(t/2\pi e)/\log(t/2\pi)\to t$
as $t\to\infty$ and $dW'$ is the rescaled white noise.
In the limit $t\to\infty$, $\vartheta(t)/\vartheta'(t)\to t$ and
$\hat{h}(\lambda/\vartheta'(t))\to\hat{h}(0)$ in $L^2$ on compact sets
of $\lambda$, and the representation converges to
\begin{equation}
  \frac{\hat{h}(0)}{\sqrt{\vartheta'(t)}}
  \int_{-\infty}^{\infty} e^{i\lambda t}\,dW'(\lambda),
\end{equation}
recovering the stationary process structure in the $t$-coordinate.

\medskip\noindent\textbf{Step 7: The explicit bridge to $\zeta$.}
To show that the final expression \eqref{eq:Zfull} equals $e^{i\vartheta(t)}\zeta(\tfrac{1}{2}+it)$, we establish the bridge from the stationary core $S$ back to $\eta(\tfrac{1}{2}+it)$ and then to $\zeta(\tfrac{1}{2}+it)$.

Recall that $\hat{h}(\omega)\,dW(\omega) = d\widetilde{Z}(\omega)$ is the spectral measure of the stationary core $S$, whose covariance is $\eta(\tfrac{1}{2}+i\cdot)$. The Fourier integral of $S$ evaluated at $\vartheta(t)$ is
\begin{equation}
  S(\vartheta(t)) = \int_{-\infty}^{\infty} e^{i\omega\vartheta(t)}\,d\widetilde{Z}(\omega) = \int_{-\infty}^{\infty} e^{i\omega\vartheta(t)}\hat{h}(\omega)\,dW(\omega).
\end{equation}
To connect this Fourier integral back to $\eta(\tfrac{1}{2}+it)$, we equate the modulated process $Z(t) = \sqrt{\vartheta'(t)}S(\vartheta(t))$ to the analytic sample path \eqref{eq:Z-via-eta} and isolate $S(\vartheta(t))$:
\begin{equation}
  \label{eq:bridge-S-eta}
  S(\vartheta(t)) = \frac{e^{i\vartheta(t)}}{\sqrt{\vartheta'(t)}\left(1-2^{1/2-it}\right)} \eta\!\left(\tfrac{1}{2}+it\right).
\end{equation}
We then connect this to $\zeta(\tfrac{1}{2}+it)$ by applying the relation $\eta(\tfrac{1}{2}+it) = \left(1-2^{1/2-it}\right)\zeta(\tfrac{1}{2}+it)$:
\begin{equation}
  \label{eq:bridge-S-zeta}
  S(\vartheta(t)) = \frac{e^{i\vartheta(t)}}{\sqrt{\vartheta'(t)}} \zeta\!\left(\tfrac{1}{2}+it\right).
\end{equation}
Multiplying both sides of \eqref{eq:bridge-S-zeta} by $\sqrt{\vartheta'(t)}$ yields the final equality:
\begin{equation}
  \label{eq:bridge-full}
  Z(t) = \sqrt{\vartheta'(t)}\int_{-\infty}^{\infty} e^{i\omega\vartheta(t)}\hat{h}(\omega)\,dW(\omega) = e^{i\vartheta(t)}\zeta\!\left(\tfrac{1}{2}+it\right).
\end{equation}
This confirms that evaluating the Fourier integral of the stationary core at $\vartheta(t)$ rigorously bridges back to $\eta$ and then to $\zeta$.
\end{proof}

\section{Recapitulation: The Four Layers}

\noindent\textbf{1. White noise:}
\[
  \mathbb{E}|dW(\omega)|^2=d\omega.
\]

\noindent\textbf{2. Spectral coloring:}
\[
  d\widetilde{Z}(\omega)=\hat{h}(\omega)\,dW(\omega),
  \qquad
  \mathbb{E}|d\widetilde{Z}(\omega)|^2=K_{1/2}(\omega)\,d\omega.
\]

\noindent\textbf{3. Stationary core in zero-rate coordinate $\tau$:}
\[
  S(\tau)=\int_{-\infty}^{\infty} e^{i\omega\tau}\,d\widetilde{Z}(\omega),
  \qquad
  \mathbb{E}[S(\tau)\overline{S(\sigma)}]=\eta\!\left(\tfrac{1}{2}+i(\tau-\sigma)\right).
\]

\noindent\textbf{4. Uniform modulation and pullback:}
\[
  \widetilde{Z}(\tau)=\sqrt{\vartheta'(\vartheta^{-1}(\tau))}\,S(\tau),
  \qquad
  Z(t)=\widetilde{Z}(\vartheta(t))
      =\sqrt{\vartheta'(t)}\,S(\vartheta(t)).
\]

\noindent\textbf{Expanding layer 4 via layers 2--3:}
\[
  Z(t)=\sqrt{\vartheta'(t)}\int_{-\infty}^{\infty}
        e^{i\omega\vartheta(t)}\hat{h}(\omega)\,dW(\omega).
\]

\noindent\textbf{Sample-path identity (analytic):}
\[
  Z(t)=e^{i\vartheta(t)}\zeta\!\left(\tfrac{1}{2}+it\right)
      =\frac{e^{i\vartheta(t)}}{1-2^{1/2-it}}
       \int_{-\infty}^{\infty} e^{i\omega t}K_{1/2}(\omega)\,d\omega.
\]

\noindent\textbf{Sample-path bridge to $\zeta$:}
\[
  S(\vartheta(t)) = \frac{e^{i\vartheta(t)}}{\sqrt{\vartheta'(t)}\left(1-2^{1/2-it}\right)}\eta\!\left(\tfrac{1}{2}+it\right) = \frac{e^{i\vartheta(t)}}{\sqrt{\vartheta'(t)}}\zeta\!\left(\tfrac{1}{2}+it\right)
\]
yielding the exact sample-path equality $Z(t) = \sqrt{\vartheta'(t)}S(\vartheta(t)) = e^{i\vartheta(t)}\zeta(\tfrac{1}{2}+it)$.

\bigskip
\noindent The stochastic representation captures the same second-order structure:
\[
  \mathbb{E}[Z(t)\overline{Z(s)}]
  =\sqrt{\vartheta'(t)\vartheta'(s)}\,
   \eta\!\left(\tfrac{1}{2}+i(\vartheta(t)-\vartheta(s))\right),
\]
which is the covariance of the uniformly modulated Gaussian process whose
stationary core has spectral density $K_{1/2}$.

\end{document}
